\abstract{Clinical documentation imposes a substantial burden on primary care physicians, with patient--clinician consultations generating diagnostically rich spoken dialogue that remains largely inaccessible to automated health information systems in its raw form. This paper presents the design and evaluation of an automated six-step pipeline that transforms unstructured clinical audio recordings directly into validated HL7 FHIR Release 4B-compliant data structures, without requiring custom model development or manual post-editing. Applied to the PriMock57 dataset of 57 mock primary care consultations, the pipeline integrates: acoustic pre-processing with stationary noise suppression and voice activity detection; speaker diarization constrained to the doctor--patient dyad; automatic speech recognition using Medical-Whisper-Large-v3, a domain-adapted Whisper model achieving a word error rate of 0.19 on PriMock57; large language model-driven transcript cleaning and physician/patient role assignment; joint clinical information extraction and FHIR resource generation via structured LLM output constrained to Pydantic schemas and assembled programmatically through the \texttt{fhir.resources} library; and structural FHIR validation with per-resource provenance tracking. A key architectural contribution is the consolidation of traditional multi-stage NLP pipelines into a single LLM extraction pass followed by schema-constrained programmatic FHIR construction, reducing hallucination risk while preserving end-to-end integration. The pipeline produces FHIR R4B Bundles containing Encounter, Condition, MedicationStatement, Observation, and Procedure resources with segment-level provenance links to source transcript segments. To our knowledge, this is the first system evaluated end-to-end from raw clinical audio to conformant FHIR resources on a public benchmark, addressing a gap in the literature where prior work treats transcription, clinical NLP, and FHIR generation as isolated problems.}
